\chapter{Background}\label{ch:literature-review}

This chapter reviews the five literatures this thesis draws on, in
dependency order: evaluation benchmarks and arena platforms
(Section~\ref{sec:benchmarks}); unsupervised taxonomy induction
(Section~\ref{sec:taxonomy-construction}); the reliability of LLM
judges (Section~\ref{sec:llm-as-a-judge}); pairwise ranking models
(Section~\ref{sec:ranking-models}); and query routing, the downstream
task that per-domain capability profiles serve
(Section~\ref{sec:routing}). Each section reports what the work
establishes and what its authors and successors found it does not
settle. The open questions that accumulate are the ones the design in
Chapter~\ref{ch:methodology} was built around, and
Section~\ref{sec:synthesis} draws them together.


\section{LLM Evaluation Benchmarks}\label{sec:benchmarks}

\subsection*{Static Benchmarks and Their Limitations}

The first generation of Large Language Model (LLM) evaluation relied
on static instruments such as \textit{Massive Multitask Language
Understanding} (MMLU)~\parencite{hendrycks2021measuring} and the
\textit{Holistic Evaluation of Language Models} (HELM)
framework~\parencite{liang2022holistic}. MMLU probes professional and
academic competence with over 15{,}000 multiple-choice questions
across 57 subjects; at release, the best-performing model (GPT-3,
175B) reached 43.9\% accuracy against an estimated human-expert
ceiling of 89.8\%~\parencite{hendrycks2021measuring}. HELM evaluated
30 models across 42 scenarios, of which 16 are core scenarios
measured on all seven metrics --- accuracy, calibration, robustness,
fairness, bias, toxicity, and efficiency --- wherever the pairing is
applicable, and was the first large-scale effort to resist reducing
model quality to a single number~\parencite{liang2022holistic}.

In both frameworks the evaluation corpus and its categorical
organisation are fixed at construction time, and three failure modes
follow from that. The first is \textit{saturation}: frontier models
now score at or near MMLU's estimated human-expert ceiling, so the
benchmark can no longer separate them~\parencite{hendrycks2021measuring,
wang2024mmlupro}. The second is \textit{contamination}: benchmark
questions circulate online and enter training corpora, and ChatGPT
and GPT-4 guess missing MMLU answer options at rates of 52\% and
57\%, far above what genuine reasoning would
predict~\parencite{deng2024contamination}. The third is
\textit{categorical brittleness}: editorially defined subject
categories are fixed at the granularity their authors chose, so they
cannot adapt as models improve and they leave blind spots wherever a
question sits between two categories. HELM's 2022-vintage scenarios
inherit the same limitation~\parencite{liang2022holistic}.

\subsection*{From MMLU to MMLU-Pro}

Saturation and contamination motivated a line of progressively harder
benchmarks, beginning with \textit{BIG-Bench Hard}'s extraction of
the 23 BIG-Bench tasks on which prior models had not outperformed
average human raters~\parencite{suzgun2022bigbench}.
\textit{MMLU-Pro}~\parencite{wang2024mmlupro} applies the same logic
to the MMLU lineage: ten answer options per question instead of four,
trivial and noisy items removed, and reasoning-focused questions
added from academic exams and textbooks across 14 domains, for
12{,}032 curated items. Accuracy drops by 16\% to 33\% relative to
MMLU and prompt sensitivity falls from 4--5\% to
2\%~\parencite{wang2024mmlupro}.

Harder questions do not address contamination, which persists in any
fixed pool. \textit{LiveBench}~\parencite{white2024livebench} refreshes
its questions on a rolling basis from sources published after common
training cutoffs, and organises them along a different axis: eighteen
tasks grouped into six skill categories --- math, coding, reasoning,
language, instruction following, and data analysis --- rather than
academic subjects. The two corpora therefore carry different
\textit{kinds} of taxonomy, and a per-domain profile built over one
does not transfer to the other. What they share is the third failure
mode, categorical brittleness, which no benchmark in this line
addresses. The categories
remain what their editors chose, and MMLU-Pro's 14 domains are the
finest published cut of that corpus.

\subsection*{Pairwise Arena Evaluation Platforms}

A parallel line supplies the signal static benchmarks lack: human
preference on open-ended prompts. \textit{Chatbot
Arena}~\parencite{chiang2024chatbotarena} is the canonical instance.
Users submit a free-form prompt, receive responses from two anonymous
models, and vote for the preferred one; by early 2024 the platform
had collected over 240{,}000 votes across more than 50 models. Votes are
aggregated with a Bradley--Terry model, confidence intervals come
from sandwich standard errors --- a variance estimate that stays valid
when the model is misspecified --- and an adaptive sampling rule
allocates comparison budget where it most reduces interval size,
reaching a given precision with fewer comparisons than uniform random
sampling~\parencite{chiang2024chatbotarena}. The
design settled several questions at once: that pairwise votes can be
gathered at scale, that Bradley--Terry is a workable aggregator for
them, and that adaptive scheduling pays for itself.

It also has a property that follows from its own premise. Because the
prompts are free-form and user-supplied, no answer key exists, so a
vote is not an estimate of anything external --- it \textit{is} the
measurement. That makes the platform's own operating conditions part
of the instrument. \textcite{singh2025leaderboard} document that
private testing of model variants before public listing, asymmetric
access to vote data, and unequal sampling and deprecation rates
across providers advantage a small set of participants, so observed
Arena rank partly reflects platform engagement rather than capability
alone. There is no external quantity against which such an effect
could be detected from the leaderboard itself.

A single global ranking also conceals variation across topics, and
the Arena line has responded by inducing structure from its own query
stream. \textit{BenchBuilder}~\parencite{li2024benchbuilder} embeds
roughly 200{,}000 Arena prompts, clusters them hierarchically into some
4{,}000 topic clusters, scores prompts with an LLM against seven
quality criteria, and distils a 500-prompt benchmark
(Arena-Hard-Auto) whose model ranking reaches 98.6\% confidence
agreement and 96.7\% Spearman correlation with Chatbot Arena's
human-vote hard-prompt leaderboard, at roughly \$20 of judging cost
per evaluated model. Arena has since induced categories directly for a
public leaderboard: the seven web-development categories on the Code
Arena leaderboard were derived from over 250{,}000 collected prompts by
clustering analysis followed by an iterative taxonomy refinement
against explicit criteria~\parencite{arena2025categories}.
BenchBuilder's clusters are
machinery for selecting prompts rather than an object of study:
membership is hard, cluster coherence is not measured, and no
per-cluster model statistics survive into the published benchmark.
Nothing in the method requires them to, since the deliverable is the
distilled prompt set.

\textit{Prompt-to-Leaderboard} (P2L)~\parencite{frick2025p2l} takes
per-topic resolution to its continuous limit. A language model
trained on Arena's human preference votes maps an individual prompt
to a vector of prompt-conditional Bradley--Terry coefficients, which
yields a leaderboard, and hence a routing decision, per prompt; a
P2L-based router reached first place on Chatbot Arena in January
2025. The capability profile is amortised into regression weights, so
the resolution is as fine as the prompt itself and the cost of a new
profile is a forward pass. The representation carries the
corresponding property: the domains are implicit, and what the model
has learned about any region of prompt space is legible only through
its outputs on prompts drawn from it.


\section{Unsupervised Taxonomy Induction}\label{sec:taxonomy-construction}

\subsection*{Corpus-Induced Taxonomies in Deployment}

The embed--cluster--LLM-label recipe has become an industrial
pattern. \textit{TnT-LLM}~\parencite{wan2024tntllm} has an LLM
generate and iteratively refine a label taxonomy over conversation
logs, then distils it into classifiers deployed over Bing Copilot
traffic. \textit{Clio}~\parencite{tamkin2024clio} extracts
privacy-preserving facets --- short model-written summaries and
attributes --- from millions of Claude conversations and embeds and
clusters those rather than the verbatim transcripts, has an LLM name
each cluster, and assembles the names into a navigable hierarchy for
usage monitoring.
\textit{TopicGPT}~\parencite{pham2024topicgpt} replaces the
clustering step itself with prompted topic generation and assignment.

These systems demonstrate that readable structure can be induced from
a large corpus and put to work. They also share a property that
follows from their purpose: the taxonomy is consumed downstream ---
for classification, monitoring, or curation --- and is validated, when
it is validated, by how well the downstream task performs.
Coherence of an individual cell is not itself a reported quantity,
and none of these systems is built to answer whether a particular
cell is a real region of the corpus or an artefact of where the
clustering happened to cut.

\subsection*{Spherical Embeddings and vMF Clustering}

When text is encoded by a neural model and $\ell_2$-normalised to the
unit hypersphere, cosine similarity becomes the natural inner product
and the von Mises--Fisher (vMF) distribution becomes the canonical
probability model for directional data~\parencite{banerjee2005vmf}.
The vMF density for a unit vector $\mathbf{x} \in \mathcal{S}^{d-1}$,
with mean direction $\boldsymbol{\mu}$ and concentration
$\kappa \geq 0$, is
\begin{equation}
    p(\mathbf{x};\boldsymbol{\mu},\kappa) =
    C_d(\kappa)\exp\!\bigl(\kappa\,\boldsymbol{\mu}^\top\mathbf{x}\bigr),
    \label{eq:vmf}
\end{equation}
where $C_d(\kappa)$ is a normalising constant involving a modified
Bessel function of the first kind. The vMF density and the EM
procedure for fitting a mixture of them are given
by~\textcite{banerjee2005vmf}; the companion paper on Bregman
clustering~\parencite{banerjee2005bregman} establishes separately
that clustering with any regular exponential family is exactly
minimisation of the corresponding Bregman divergence, which is what
places vMF mixture fitting in that framework. High-order
approximations to the concentration estimate remain numerically
stable in the high-$\kappa$, high-dimensional regime of neural
embedding spaces~\parencite{banerjee2005vmf, hornik2014vmf}.
Because the model is generative, a candidate split has a likelihood
and the improvement it buys is a quantity rather than a threshold
crossing, which is what makes an acceptance statistic available at
all.

A recursive construction over queries also needs an embedding space
in which coarse domain-level and fine sub-topic distinctions are both
resolvable, and a single fixed-width representation is a compromise
between them. Matryoshka Representation
Learning~\parencite{kusupati2022mrl} removes the compromise by
optimising one embedding jointly across nested prefix dimensions, so
the most discriminative directions concentrate in the leading
coordinates and any prefix is a coherent representation in its own
right. A fixed prefix width can therefore be used at every depth
without placing nodes at different depths on different spheres, which
is the condition under which cross-depth comparisons mean anything.

\subsection*{Overlapping and Polyhierarchical Structure}

Relaxing the topology so that a node may have more than one parent
yields a \textit{polyhierarchy}, a directed acyclic
graph~\parencite{polyhierarchy2006} --- standard in knowledge
organisation, where it is curated by hand, but rare in corpus-induced
construction, with correspondingly little evidence about whether
learned multi-parent structure earns its additional complexity.


\section{LLM-as-a-Judge and Reliability}\label{sec:llm-as-a-judge}

As human evaluation becomes a scalability bottleneck, LLM-as-a-judge
has become the dominant approach to automated assessment of model
outputs~\parencite{zheng2023judging}: a frontier model is asked to
evaluate the semantic quality of candidate responses. Its reliability
properties bound any framework that delegates comparison to it.

\subsection*{Pointwise and Pairwise Judging}

Pointwise evaluation asks a judge model to score a single response on
an absolute scale; pairwise evaluation presents two responses and
asks for a preference~\parencite{zheng2023judging}. Benchmarking the
two framings against human preference with MT-Bench and Chatbot
Arena, Zheng et al.\ report agreement with human raters at or above
the level of human--human agreement for strong judges, and note that
the relative task sidesteps the absolute score calibration a
pointwise judge has to supply. Their evidence bears more directly on
agreement with humans than on sensitivity to prompt formatting, and
the comparison between the two framings is drawn across settings
rather than from a controlled head-to-head. It is nonetheless the
main empirical basis for the pairwise framing that arena evaluation
is built on.

\subsection*{Systematic Biases}

Three systematic biases distort pairwise preferences independently of
response quality.

\begin{itemize}
    \item \textbf{Positional bias.} Judges tend to favour the
    response presented first in the prompt, independently of
    content~\parencite{zheng2023judging}. Zheng et al.\ report a
    positional consistency of 65.0\% for GPT-4: in only about two
    thirds of pairs did the judge reach the same verdict when the
    positions of the two responses were swapped, so about a third of
    single-order outcomes are order-sensitive. Inconsistency is not
    the same as position-driven, since noise also contributes, but a
    third is a first-order term.

    \item \textbf{Verbosity bias.} Models assign higher scores to
    longer responses even when the added content contributes no
    reasoning depth. \textcite{zheng2023judging} first flag the
    effect, and \textcite{saito2023verbosity} characterise it
    directly in preference labelling, reporting a gap between the
    length preference of LLM labellers and that of human ones. The
    confound is
    acute in knowledge-intensive question answering, where a model
    can pad a terse correct answer with confident elaboration and win
    a preference it should not have won.

    \item \textbf{Self-preference bias.} LLMs favour outputs
    stylistically similar to their own training
    distribution~\parencite{wataoka2025selfpreference,
    panickssery2024llm}, and \textcite{panickssery2024llm} link the
    effect to self-recognition: evaluators that can identify their
    own generations are the ones that favour them.
\end{itemize}

The three differ in how far a prompt can reach them. Positional bias
is the most tractable, since presenting both orders and checking
consistency identifies order-driven verdicts directly and costs one
extra call. Verbosity bias is partly addressable through explicit
evaluation criteria, with effectiveness an empirical question per
domain, and is addressable post hoc by modelling length as a
covariate and reporting length-controlled win
rates~\parencite{dubois2024lengthcontrolled}. Self-preference is
documented as a property of the evaluator rather than of the
prompt~\parencite{panickssery2024llm}; on that reading --- an
inference from the evidence rather than a finding any of these papers
states --- no instruction should be expected to remove it, and the
remedy available is to change or rotate the judge.

\subsection*{Judges on Verifiable Tasks: The Shortcut Result}

The biases above concern open-ended comparison, where no answer key
exists. On tasks that have one, a further and more consequential
behaviour is documented. Examining LLM judges on mathematical
reasoning, \textcite{stephan2024calculation} find that judges tend to
award the win to the generally stronger model even on instances where
that model's answer is wrong; that a large fraction of judgments
(70--75\%) can be predicted from shallow features without consulting
the mathematics at all; and that judges are able to detect the
stratum in which both candidate answers are correct, precisely the
stratum where a preference verdict is least informative. Verdicts, in
short, track answer-level correctness and prior model strength rather
than the quality of the reasoning being adjudicated.

Meta-evaluation benchmarks corroborate the pattern from the other
direction. \textit{JudgeBench}~\parencite{tan2025judgebench}
constructs response pairs whose preference label is objective
correctness and finds that the judges it evaluates --- including
strong general-purpose models and dedicated reward models --- perform
near chance on them, a result reported for that set of judges at
publication rather than as a ceiling, since later
reasoning-optimised judges have scored well above it on the
same split; \textit{LLMBar}~\parencite{zeng2024llmbar} shows
judges systematically fail pairs in which superficial quality and
instruction adherence are dissociated; and
\textit{RewardBench}~\parencite{lambert2024rewardbench} standardises
this style of adversarial audit for reward models. Judge competence,
these benchmarks agree, does not transfer as a scalar from fluent
domains to verifiable ones.

The evidence is largely post-hoc: a corpus of verdicts is collected
and then analysed for what predicts it. That establishes the
phenomenon but leaves the mechanism underdetermined, because a judge
that verifies the answer and a judge that applies quality criteria
agree on most instances and are separated only on the strata where
the criteria have nothing to work with --- both responses correct,
neither correct, no reasoning present to evaluate. Those strata exist
in any keyed multiple-choice corpus and can be identified before any
verdict is collected.

\subsection*{Domain Calibration}

A judge prompted only to ``prefer the better answer'' applies
implicit quality criteria drawn from its training distribution, which
may not match the standards of the domain.
\textit{FLASK}~\parencite{ye2024flask} makes this concrete: when a
judge scores each instruction against the specific skills that
instruction requires, rather than against a single holistic
preference, the relative ordering of models differs from one skill to
the next and agreement with human raters improves, so the choice of
criteria is part of the measurement, not a presentational detail.

The mitigations differ in what supervision they consume. The earliest
is Zheng et al.'s reference-guided mode, in which the judge is shown
a reference solution at verdict time~\parencite{zheng2023judging}.
\textit{MemAlign}, a non-peer-reviewed system description cited as an
industrial instance only~\parencite{databricks2026memalign}, aligns
synthetic judges with domain standards through a dual memory: a
semantic store of general evaluation principles distilled from expert
feedback, and an episodic store of task-specific cases in natural
language, retrieved together at verdict time. Its reported training
sets are on the order of 50 human gold labels, with improvement
claimed from far fewer; those figures are vendor-reported and have
not been through peer review. On the training side, the RL-trained
reasoning judges
\textit{Think-J}~\parencite{huang2025thinkj} and
\textit{JudgeLRM}~\parencite{chen2025judgelrm} reward explicit
chain-of-thought justifications consistent with the final preference,
which makes the evaluation auditable and reduces preference reversals
attributable to surface features. Domain conditioning improves
alignment with human expert ratings across this family, at the cost
of labelled examples or a trained reasoning judge.

One arrangement is absent from it. Reference-guided judging reveals
the reference when the verdict is made, and FLASK instantiates its
criteria per instance; both keep the domain signal and the judged
instance together. Distilling criteria once from a domain's answer
key and then withholding that key at verdict time separates the two,
and how a judge behaves under that separation --- in particular
whether the distilled criteria do any work once the key is gone --- is
not something the calibration literature reports.


\section{Pairwise Ranking Models}\label{sec:ranking-models}

\subsection*{Elo and Its Limitations}

The \textit{Elo rating system}~\parencite{elo1978rating} fits LLM
evaluation poorly: it assumes static skill between matches and models
no uncertainty, so rankings become unstable under sparse interaction
data. \textcite{boubdir2024elo} find that for models with win rates
near 50\% the ordering is sensitive to the $K$-factor and to the
number of match permutations averaged over --- an instability of the
rating system under near-tied comparisons, not a measurement of
cyclic preference among models. That instability among near-tied
models is what motivates a likelihood-based aggregate.

\subsection*{Bradley--Terry and Bayesian Extensions}

The \textit{Bradley--Terry} (BT) model provides the formal
probabilistic link for pairwise preferences, modelling the merit of
each model as a latent parameter
$\lambda_i > 0$~\parencite{bradley1952rank}. The probability that
model $i$ is preferred over model $j$ is
\begin{equation}
    P(i \succ j) = \frac{\lambda_i}{\lambda_i + \lambda_j}
    = \sigma(\theta_i - \theta_j),
    \label{eq:bt}
\end{equation}
where $\theta_i = \log \lambda_i$ and $\sigma$ is the logistic
function. Bradley and Terry give the model and its maximum-likelihood
estimator; the surrounding computational and asymptotic theory was
formalised later. The log-likelihood is strictly concave in
$(\theta_i)$ and the MLE can be obtained by iterative re-weighted
least squares or by the minorisation--maximisation (MM) algorithm
of~\textcite{hunter2004mm}, and \textcite{simons1999asymptotics}
establish consistency and asymptotic normality, so the asymptotic
covariance is the inverse Fisher information matrix and yields
per-parameter standard errors $\hat{\sigma}_i$. Those standard errors
are what make the model more than a scoring rule: they are consumed
downstream by any adaptive scheduler and by any stopping criterion.

Two conditions bound when the estimate exists and means what it
appears to, and both are due to \textcite{ford1957ranking} rather
than to the original paper. Identifiability requires a connected
comparison graph (models as nodes, an edge per observed comparison);
if the graph decomposes, BT parameters are identified only up to an
additive constant per component, so scores from different components
are not on a common scale and a difference between them is not a
quantity. Connectivity is necessary but not sufficient: Ford's
condition is that the graph of \emph{outcomes} be strongly connected,
and complete separation --- one model winning every comparison against
another --- violates it, sending the unregularised MLE to infinity.
The same strong-connectivity requirement appears as Assumption~1
of~\textcite{hunter2004mm}, where it is what guarantees the MM
iteration converges. Both conditions bite hardest where
comparisons are few, which is exactly the regime a per-domain arena
creates when it divides a fixed budget across many cells.

Bayesian alternatives exist at both ends: online-update systems for
streaming match data (TrueSkill~\parencite{herbrich2007trueskill},
OpenSkill~\parencite{weng2011bayesian}), and hierarchical
partial-pooling BT models that shrink per-group strengths toward a
population mean~\parencite{cattelan2012paired} --- stabilising
sparse groups at the cost of making each group's estimate depend on
the others, a trade that matters when the groups are themselves the
object of study.

\subsection*{Active Ranking and Sample Complexity}

Round-robin evaluation of $M$ models requires $\binom{M}{2}$
comparison pairs per cell. \textit{Active ranking} reduces this by
selecting, at each step, the comparison most likely to reduce global
ranking uncertainty; \textcite{jamieson2011activeranking} established
that adaptive selection can recover a ranking in far fewer
comparisons than the pair count suggests.

In the stochastic setting, where comparison outcomes are noisy
Bernoulli variables as with an LLM judge,
\textcite{heckel2019activeranking} derive finite-sample bounds under
the BT model directly: distinguishing two adjacent models with error
probability $\delta$ requires $O(\Delta^{-2} \log(M/\delta))$
comparisons, where $\Delta$ is the minimum separation between
consecutive BT parameters. Close pairs need disproportionately many
comparisons, so uniform budget allocation is inefficient, and a rule
that schedules the pair with the largest overlap in BT confidence
intervals achieves the bound without knowledge of
$\Delta$~\parencite{heckel2019activeranking}. Heckel et al.\ further
establish that a parametric model (BTL, Thurstone) offers at most
$O(\log n)$ improvement over non-parametric adaptive strategies, so
the case for BT rests on interpretability and on the convergence
criterion its standard errors provide rather than on sample
efficiency alone. A scheduler of this family reads those standard
errors, which makes its behaviour inseparable from the correctness of
the variance it is given.

\section{Query Routing as the Target Use Case}\label{sec:routing}

\subsection*{The Routing Problem in Agentic Systems}

In multi-model agentic systems, a query is dispatched to the
candidate model expected to perform best on it, which requires
knowledge of each model's domain-conditional performance. If model
$A$ outperforms model $B$ on mathematical reasoning but underperforms
on legal analysis, a global ranking that averages across domains
mis-routes queries at both extremes.
\textcite{raju2024constructing} supply the evaluation-side companion
to that premise, though a narrower one than is sometimes read into
them: scoring ten top-ranked models over a curated set of fourteen
domains, they find that per-category \emph{separability} --- the
fraction of model pairs whose confidence intervals do not overlap ---
varies substantially by domain, and they attribute the variation to
per-category sample size rather than to models reordering. Whether
per-domain \emph{rankings} genuinely differ is not a quantity their
pipeline measures. Their study measures the evaluation sets rather
than any router, and it presupposes what it uses: per-domain scores
that already exist, over a domain taxonomy that is given or
hand-built.

\subsection*{Learned Routing Systems}

\textit{FrugalGPT}~\parencite{chen2023frugalgpt} cascades through
models of increasing capability until a learned confidence estimator
accepts the current response, and reports inference-cost reductions
of up to 98\% at matched accuracy.
\textit{RouteLLM}~\parencite{ong2024routellm} trains binary routing
functions, among them a similarity-weighted Bradley--Terry router, on
80{,}000 pairwise preference judgements from Chatbot Arena. Its most
informative negative result is about the training signal: routers
trained solely on the global preference distribution generalise
poorly to domain-specific benchmarks such as MMLU and GSM8K unless
the training data is augmented with in-domain labels, because a
globally trained BT model does not encode the domain-conditional
variation accurate routing requires.
\textit{EmbedLLM}~\parencite{zhuang2024embedllm} and
\textit{UniRoute}~\parencite{jitkrittum2025uniroute} route
geometrically, representing each model by an embedding derived from
its performance on a representative query set. Across all of these,
the domain structure over which models are profiled is either absent
or inherited from whichever benchmark supplied the training set.

A more recent line induces its routing structure directly from the
query corpus --- embed-and-cluster routers that score each candidate
model per cluster
(\textit{The Avengers}~\parencite{zhang2025avengers,
zhang2025avengerspro}), error-profile and jointly learned embeddings
(UniRoute's cluster variant~\parencite{jitkrittum2025uniroute};
\textit{RouterDC}~\parencite{chen2024routerdc}), with
\textit{RouterBench}~\parencite{hu2024routerbench} standardising
their comparison.

Their structure is tuned against routing accuracy, which is the
appropriate objective for a router and has a consequence for what the
structure can be said to be. A partition selected because it routes
well is evidence about the routing pipeline as a whole; it is not
independent evidence that the partition captures real regions of the
corpus, since the two are optimised together and cannot be
attributed separately after the fact. Establishing the second
requires validating the partition against something other than the
task it was built to serve.


\section{Synthesis}\label{sec:synthesis}

Read together, the five literatures leave a specific arrangement
untried rather than a large gap unfilled.

Every component of the arrangement below already exists on its own.
Corpus-induced structure over an evaluation query stream is a solved
engineering problem and an industrial pattern (BenchBuilder, TnT-LLM,
Clio). Per-topic and per-prompt capability profiles exist: Arena
publishes per-category leaderboards~\parencite{arena2025categories,
arena2025expert},
and P2L~\parencite{frick2025p2l} carries the idea to a
prompt-conditional Bradley--Terry vector. Uncertainty quantification
for task-specific rankings under sparse comparison budgets is itself a
subject of study, as low-rank estimation~\parencite{li2026lowrank}, as
semiparametric tensor completion~\parencite{li2026tensor}, and as
conformal rank intervals over a
leaderboard~\parencite{neuhof2026rankintervals}; each of these takes
the task partition as given. Intrinsic
validation of cluster coherence, independent of any downstream task,
exists as well --- \textcite{hong2025dialin} fit an LLM to judge
semantic coherence and iterate cluster discovery against it, and
\textcite{kargupta2025taxoadapt} report coherence as a headline
metric of an induced taxonomy. Corpus-induced routing structures
exist (The Avengers, UniRoute's cluster mode, RouterDC). Bradley--Terry
supplies a pairwise model with per-parameter standard errors, and
active ranking supplies a scheduling rule that consumes them.

What no line couples is intrinsic coherence validation with per-cell
estimation. Where coherence is validated
intrinsically, nothing is estimated per cell; where per-cell or
per-prompt Bradley--Terry coefficients are estimated, the cell is
taken as given --- inherited from an editorial taxonomy, selected
because it routes well, or, in P2L's case, not a discrete object that
could be audited one at a time.

The nearest case is
\textcite{maimon2025skills}, who recover latent skills by factor
analysis over a model-by-benchmark matrix and report each factor's
internal consistency (McDonald's $\omega > 0.80$) before relating
those skills to Arena scores; that coherence is the psychometric
reliability of a latent factor rather than a test that a region of a
query corpus is real, and the factors are continuous latent dimensions
rather than discrete auditable cells with a pass/fail gate preceding
estimation.

The novelty claimed here is that coupling of validation and
estimation, and nothing wider: a per-cell coherence and acceptance
gate that a cell must pass \textit{before} a Bradley--Terry profile is
estimated inside it, so that the estimate and the region it describes
are certified by the same construction.

Building the partition as an instrument in its own right --- discrete
cells with a tested coherence criterion on the query side and
Bradley--Terry profiles with standard errors on the model side,
validated before any router consumes them --- is the arrangement this
thesis constructs and measures. Chapter~\ref{ch:methodology}
specifies it and Appendix~\ref{app:dag-formalism} derives its formal
foundations.

The corpus decides what such a measurement can establish, and in one
respect it cuts both ways. Because MMLU-Pro carries a verifiable
answer key, each link in the chain can be checked against ground
truth rather than against human preference, which is the property
that makes a keyed corpus worth using at all. Because it carries a
verifiable answer key, the shortcut literature of
Section~\ref{sec:llm-as-a-judge} predicts that a capable judge may
reach its verdict by checking the answer rather than by applying the
criteria it was given. The same property that makes the measurement
possible is the one that bounds what a rubric can be shown to do.
Whether that prediction holds here is the subject of
Chapter~\ref{ch:results}; this chapter states the expectations
against which those results are read, and does not anticipate them.
